\subsection{A training-selected vector-fetch path}
\label{sec:profile-protocol}
The scheduling reversals suggest a modest execution proposal: choose how to fetch an expansion's vectors using a short profile of the target configuration. HNSW still chooses which candidates to inspect. The optional path chooses only the order in which their memory transactions are offered. It requires a near-memory engine or frontend that exposes those requests; the experiment does not establish that a particular commercial PIM product provides this interface.

For each reference split, use the first 50 of its 500 training queries. Replay their BFS vector placement with original order, address sorting and transaction round-robin under the target configuration. Select the order with minimum total completed service cycles; exact ties prefer original order, then sorting, then round-robin. Store that choice once and use it for all 200 held-out queries. This is a configuration-level choice, not a decision made using each test query's outcome. The training and test IDs are disjoint within every split.

The profile is specific to the graph/layout, footprint, batch limit, outstanding bound and memory configuration. It selects only request order; it does not choose a layout, batch size or hardware configuration. The experiment repeats it for the default and six sensitivity conditions at both widths, producing 140 choices from 420 calibration replays. The selected order's held-out request stream has already been measured in the scheduling campaigns; evaluation uses that exact record. Calibration runs are additional training evidence, not additional independent performance repetitions.

This extension was designed after inspecting early sensitivity reversals and previously evaluated test outcomes. Its 50-query budget, three choices and tie rule were then fixed in \texttt{doc/PROFILED\_FETCH\_DESIGN.md}; all outcomes are retained. Although selection code reads training costs only, the overall study remains exploratory because those observations motivated the policy and queries recur across splits. An untouched workload would be needed for a stronger generalisation claim.

Compare the selected path with every fixed schedule under matching settings. Also report its difference from the retrospectively best fixed schedule on each held-out split. That \emph{oracle} is a descriptive reference: it uses test outcomes and is unavailable to a deployed selector. At default settings, compare with the measured RCM, Gorder and Porder placements using original order. Those are whole-configuration comparisons; combining published placements with the selector is a separate, unevaluated intervention. Report profiling service cost separately from held-out savings. No inference about application amortisation follows without implementation costs in common units.
